
\subsubsection*{1.3.6}
\phantomsection
\label{sga3:VIIB:1.3.6}
\addcontentsline{toc}{subsubsection}{\SGARefLink{sga3:VIIB:1.3.6}{1.3.6}. Filtered flat coalgebras}

In the remainder of this Expos\'e, we shall several times define a
topologically flat formal \(k\)-variety \(X\) by exhibiting the
coalgebra \(\mathbf H(X)\). We shall then have to interpret, by means
of \(\mathbf H(X)\), certain geometric properties of \(X\). Here is
one example of this situation: suppose that a section \(\sigma\) of
the structure morphism \(X\to\operatorname{Spf}(k)\) is given, and
ask under what condition \(\sigma\) induces an isomorphism of the
underlying topological spaces.\footnote{\label{sga3:VIIB:nde:62}In
\SGARefLink{sga3:VIIB:1.3.6.A}{Lemma 1.3.6.A} below, the editors supplied detailed proofs of
parts \textup{\SGARefLink{sga3:vii:paragraph:5-1:equation:i}{(i)}} and \textup{\SGARefLink{sga3:vii:paragraph:5-1:equation:ii}{(ii)}} and added part \textup{\SGARefLink{sga3:VIIB:1.3.6.A:item:iii}{(iii)}}.}

\textit{To begin, suppose that \(k\) is Artinian.} Let
\((H,\Delta,\varepsilon)\) be a flat \(k\)-coalgebra, put
\(H^+=\ker(\varepsilon)\), and let \(A=H^*\) be the profinite
\(k\)-algebra dual to \(H\). Suppose that a morphism of
\(k\)-coalgebras \(k\to H\) is given; that is, an element
\(\phi\in H\) such that
\[
  \varepsilon(\phi)=1,
  \qquad
  \Delta(\phi)=\phi\otimes\phi.
\]
On the one hand, \(\phi\) defines a continuous algebra homomorphism
\(\Phi:A\to k\), and hence a section
\[
  \sigma:\operatorname{Spf}(k)\longrightarrow\operatorname{Spf}(A)
\]
of the projection
\(\operatorname{Spf}(A)\to\operatorname{Spf}(k)\).

On the other hand, define \(k\)-submodules of \(H\) by setting
\(H_0=k\phi\) and, for \(n\geq1\),
\[
  H_n
  =
  \{x\in H\mid
    \Delta(x)-x\otimes\phi\in H_{n-1}\otimes H^+\}.
\]
This formula also applies when \(n=0\) if \(H_{-1}=0\). Induction on
\(n\) shows that \(H_{n-1}\subset H_n\). We call
\[
  H_0\subset H_1\subset\cdots
\]
the \emph{filtration of \(H\) defined by \(\phi\)}.

\medskip
\noindent\textbf{Remark.}
\phantomsection
\label{sga3:VIIB:1.3.6.remark}
Since
\[
  \Delta(H_n)
  \subset
  H_n\otimes H_0\oplus H_{n-1}\otimes H^+,
\]
we have \(\Delta(H_n)\subset H_n\otimes H\). Since \(\Delta\) is
cocommutative---that is,
\(\tau\circ\Delta=\Delta\), where
\(\tau(a\otimes b)=b\otimes a\)---we likewise have
\(\Delta(H_n)\subset H\otimes H_n\). When \(H/H_n\) is flat over
\(k\), it follows that \(H_n\) is a subcoalgebra of \(H\)
\textup{(see also \SGARefLink{sga3:VIIB:1.3.6.A}{1.3.6.A}\,\textup{\SGARefLink{sga3:VIIB:1.3.6.A:item:iii}{(iii)}} below)}. In general,
however, the map
\[
  \Delta:H_n\longrightarrow H_n\otimes H
\]
does not factor through \(H_n\otimes H_n\).

\medskip
\noindent\textbf{Lemma 1.3.6.A.}
\phantomsection
\label{sga3:VIIB:1.3.6.A}
{\itshape Let \(k\) be an Artinian ring, let \(H\) be a flat
\(k\)-coalgebra, and let \(A=H^*\) be its dual profinite
\(k\)-algebra. Let \(\phi\in H\) satisfy
\(\varepsilon(\phi)=1\) and
\(\Delta(\phi)=\phi\otimes\phi\). Let
\(\Phi:A\to k\) be the continuous algebra homomorphism, let
\(\sigma:\operatorname{Spf}(k)\to\operatorname{Spf}(A)\) be the
section of
\(\operatorname{Spf}(A)\to\operatorname{Spf}(k)\), and let
\((H_n)\) be the filtration of \(H\) defined by \(\phi\). Put
\(I=\ker\Phi\).
\begin{enumerate}
  \item[\textup{(i)}]
  For every \(n\geq1\), \(H_{n-1}\) is the orthogonal in \(H\) of
  the closure \(\overline{I^n}\) of \(I^n\).

  \item[\textup{(ii)}]
  Consequently, \(\sigma\) induces a bijection of the underlying
  sets if and only if
  \[
    H=\bigcup_n H_n.
  \]

  \item[\textup{\SGAInlineRefTarget{sga3:VIIB:1.3.6.A:item:iii}(iii)}]
  If, in addition, every \(H/H_n\) is flat over \(k\), then for every
  \(n\geq0\),
  \begin{equation}
    \Delta(H_n)
    \subset
    \sum_{i=0}^n H_i\otimes H_{n-i}.
    \tag{*}
    \label{sga3:VIIB:1.3.6.A:star}
  \end{equation}
  In particular, every \(H_n\) is then a subcoalgebra of \(H\).
\end{enumerate}}

\textit{Proof.}
For every \(x\in H\), the map
\[
  A\longrightarrow k,\qquad f\longmapsto f(x),
\]
is continuous. Thus, if \(I^n\) is annihilated by \(x\), so is its
closure \(\overline{I^n}\). For \(n\geq1\), put
\[
  (\overline{I^n})^\perp
  =
  \{x\in H\mid
    f(x)=0\text{ for every }f\in\overline{I^n}\}.
\]

Suppose first that
\[
  \sigma:\mathfrak m\longmapsto\Phi^{-1}(\mathfrak m)
\]
is a bijection from \(\operatorname{Spf}(k)\) onto
\(\operatorname{Spf}(A)\). Since \(I\) is contained in the
intersection of the ideals \(\Phi^{-1}(\mathfrak m)\), it follows
from \SGARefLink{sga3:VIIB:0.1.2}{0.1.2} that the sequence of ideals
\((\overline{I^n})\) tends to zero. Let \(x\in H\). The submodule
\[
  J(x)=\{f\in A\mid f(x)=0\}
\]
is open and closed in \(A\); hence it contains
\(\overline{I^n}\) for all sufficiently large \(n\). In other words,
\(x\in(\overline{I^n})^\perp\) for all sufficiently large \(n\).

Conversely, suppose that
\[
  H=\bigcup_n(\overline{I^n})^\perp.
\]
Let \(\mathfrak p\) be an open prime ideal of \(A\). By the
definition of the topology on \(A\) \textup{(\SGARefLink{sga3:VIIB:0.2.2}{0.2.2})},
\(\mathfrak p\) contains an open \(k\)-submodule of the form
\[
  \mathscr V(x_1,\ldots,x_s)
  =
  \{f\in A\mid f(x_1)=\cdots=f(x_s)=0\}.
\]
By hypothesis, there is an integer \(n\) such that
\[
  x_1,\ldots,x_s\in(\overline{I^n})^\perp;
\]
therefore \(\overline{I^n}\subset\mathfrak p\). Moreover, since
\(k\) is Artinian,
\[
  \operatorname{Spf}(k)=\{\mathfrak m_1,\ldots,\mathfrak m_r\}
\]
is finite, and there is an integer \(t\geq1\) such that
\[
  \left(\prod_i\mathfrak m_i\right)^t=0.
\]
It follows that
\[
  \prod_i\Phi^{-1}(\mathfrak m_i)^t\subset I.
\]
Thus \(\mathfrak p\) contains the product of the ideals
\(\Phi^{-1}(\mathfrak m_i)^{tn}\). Since \(\mathfrak p\) is prime,
it contains, and hence equals, one of the
\(\Phi^{-1}(\mathfrak m_i)\). We have therefore proved that
\begin{equation}
  \sigma\text{ is a bijection}
  \quad\Longleftrightarrow\quad
  H=\bigcup_n(\overline{I^n})^\perp.
  \tag{65}
  \label{sga3:VIIB:1.3.6:eq65}
\end{equation}

We also have \(H=k\phi\oplus H^+\); let
\(\pi:H\to H^+\) be the projection with kernel \(k\phi\). For
\(n\geq0\), let
\[
  \Delta^n:H\longrightarrow H^{\otimes(n+1)}
\]
be the iterated comultiplication, let
\(\overline{\Delta}^{\,n}\) be the composite of \(\Delta^n\) with
the projection
\[
  \pi^{\otimes(n+1)}:
  H^{\otimes(n+1)}
  \longrightarrow
  (H^+)^{\otimes(n+1)},
\]
and put
\[
  H_n'
  =
  \ker(\overline{\Delta}^{\,n})
  =
  \left\{
    x\in H\ \middle|\
    \Delta^n(x)\in
    \sum_{i=0}^n
      H^{\otimes(n-i)}\otimes H_0\otimes H^{\otimes i}
  \right\}.
\]
We set \(\overline{\Delta}^{\,0}=\operatorname{id}_H\), so that
\(H_0'=H_0\). A straightforward induction on \(n\) gives
\begin{equation}
  H_n
  \subset
  H_n'
  \subset
  (\overline{I^{n+1}})^\perp.
  \tag{*}
  \label{sga3:VIIB:1.3.6.A:star-prime}
\end{equation}

Up to this point we have not used the hypothesis that \(H\) is flat
over \(k\). Assume now that \(H\) is flat, and hence projective, over
\(k\). Then \(A^\dagger=H\) by \SGARefLink{sga3:VIIB:0.2.2}{0.2.2}. We shall prove that
\[
  H_n=(\overline{I^{n+1}})^\perp.
\]
This is clear for \(n=0\). Suppose it has been proved for \(n<r\).
Then \(H_{r-1}\) is the kernel of the homomorphism
\[
  H\longrightarrow(\overline{I^r})^\dagger.
\]
Since \(H^+\) is flat, the homomorphism
\[
  (H/H_{r-1})\otimes H^+
  \longrightarrow
  (\overline{I^r})^\dagger\otimes H^+
\]
is injective. On the other hand, the hypothesis implies that \(I\)
is a projective pseudocompact \(k\)-module, since it is a direct
factor of \(A=H^*\). Thus, by \SGARefLink{sga3:VIIB:0.3.6}{0.3.6},
\[
  \bigl(\overline{I^r}\widehat{\otimes}I\bigr)^\dagger
  \simeq
  (\overline{I^r})^\dagger\otimes I^\dagger
  =
  (\overline{I^r})^\dagger\otimes H^+.
\]
The exact sequence
\[
  \overline{I^r}\widehat{\otimes}I
  \longrightarrow A
  \longrightarrow A/\overline{I^{r+1}}
  \longrightarrow0
\]
therefore gives, by duality, the exact sequence
\begin{equation}
  (\overline{I^r})^\dagger\otimes H^+
  \xleftarrow{\ \delta\ }
  H
  \longleftarrow
  (A/\overline{I^{r+1}})^\dagger
  \longleftarrow0,
  \tag{1}
  \label{sga3:VIIB:1.3.6.A:eq1}
\end{equation}
where \(\delta\) is obtained by composing \(\Delta_H\) with the
projection
\begin{equation}
  H\otimes H
  \longrightarrow
  H\otimes H^+
  \longrightarrow
  (H/H_{r-1})\otimes H^+
  \lhook\joinrel\longrightarrow
  (\overline{I^r})^\dagger\otimes H^+.
  \tag{2}
  \label{sga3:VIIB:1.3.6.A:eq2}
\end{equation}
For every \(u\in H\), the projection of \(\Delta(u)\) onto
\(H\otimes H^+\) is \(\Delta(u)-u\otimes\phi\). Equations
\textup{\SGARefLink{sga3:VIIB:1.3.6.A:eq1}{(1)}} and \textup{\SGARefLink{sga3:VIIB:1.3.6.A:eq2}{(2)}} show that if
\[
  u\in
  (A/\overline{I^{r+1}})^\dagger
  =
  (\overline{I^{r+1}})^\perp,
\]
then \(\Delta(u)-u\otimes\phi\) belongs to the kernel of
\[
  H\otimes H^+
  \longrightarrow
  (H/H_{r-1})\otimes H^+,
\]
namely \(H_{r-1}\otimes H^+\). Thus \(u\in H_r\). This completes the
proof of \textup{\SGARefLink{sga3:vii:paragraph:5-1:equation:i}{(i)}} and \textup{\SGARefLink{sga3:vii:paragraph:5-1:equation:ii}{(ii)}}, and also proves that
\[
  H_n=\ker(\overline{\Delta}^{\,n}).
\]

We now prove \textup{\SGARefLink{sga3:VIIB:1.3.6.A:item:iii}{(iii)}}. For \(i\geq0\), put
\[
  H_i^+=H_i\cap H^+.
\]
Let \(n\geq1\). For every \(x\in H_n\), the element
\(\overline x=x-\varepsilon(x)\phi\) belongs to \(H_n^+\), and
\[
  \Delta(x)
  =
  \varepsilon(x)\phi\otimes\phi
  +\overline x\otimes\phi
  +\phi\otimes\overline x
  +\overline{\Delta}(\overline x).
\]
It is therefore enough to prove that
\[
  \overline{\Delta}(H_n^+)
  \subset
  \sum_{i=1}^{n-1}H_i^+\otimes H_{n-i}^+.
\]
For \(i=0,\ldots,n-1\), the map
\(\overline{\Delta}^{\,n}\) factors as follows:
\begin{center}
  \phantomsection
  \label{sga3:VIIB:1.3.6.A:iterated-reduced-comultiplication}
  \begin{tikzcd}[row sep=3.4em,column sep=3.8em]
H^+ \arrow[r,"\overline\Delta"]
  & H^+\otimes H^+
      \arrow[r,"{\overline\Delta^{\,i}\otimes
                    \overline\Delta^{\,n-i-1}}"]
      \arrow[d]
  & (H^+)^{\otimes(i+1)}\otimes(H^+)^{\otimes(n-i)} \\
{}
  & \displaystyle
    \frac{H^+}{H_i^+}\otimes\frac{H^+}{H_{n-i-1}^+}
      \arrow[r,"f"]
  & \displaystyle
    \frac{H^+}{H_i^+}\otimes(H^+)^{\otimes(n-i)}
      \arrow[u,"g"]
      \, .
\end{tikzcd}

\end{center}
Moreover, since \(H^+/H_i^+\) and
\((H^+)^{\otimes(n-i)}\) are flat, the maps \(f\) and \(g\) in this
diagram are injective. It follows that
\[
  \overline{\Delta}(H_n^+)
  \subset
  H_i^+\otimes H^+
  +
  H^+\otimes H_{n-i-1}^+
\]
for every \(i=0,\ldots,n-1\). Part \textup{\SGARefLink{sga3:VIIB:1.3.6.A:item:iii}{(iii)}} now follows from
the next lemma, applied with \(M=H^+\), \(E_0=0\), and
\(E_i=H_{i-1}^+\) for \(i\geq1\).

\medskip
\noindent\textbf{Lemma 1.3.6.B.}
\phantomsection
\label{sga3:VIIB:1.3.6.B}
{\itshape Let \(k\) be a ring, and let
\[
  0=E_0\subset E_1\subset\cdots\subset E_n\subset M
\]
be \(k\)-modules. Suppose that \(M/E_i\) is flat for every \(i\).
Then
\footnote{The French source prints \(M\otimes E_i\), rather than
\(M\otimes E_{n-i}\), in the second summand on the left. The proof
immediately below, in particular its use of \(M/E_{n-i+1}\), requires
the latter indexing, which is used here.}
\[
  \bigcap_{i=0}^n
    \bigl(E_i\otimes M+M\otimes E_{n-i}\bigr)
  =
  \sum_{i=1}^n E_i\otimes E_{n-i+1}.
\]}

\textit{Proof.}
Denote the left-hand and right-hand sides by \(K\) and \(S\),
respectively. Clearly \(S\subset K\); we prove the reverse inclusion.
For \(i=0,\ldots,n\), put
\[
  K_i=K\cap(E_i\otimes M).
\]
For \(i=1,\ldots,n\), both \(M/E_{n-i+1}\) and
\(E_i/E_{i-1}\) are flat. Consequently, the homomorphism
\(\tau_i\) below is injective, and the composite
\[
  (E_i/E_{i-1})\otimes M
  \longrightarrow
  (E_i/E_{i-1})\otimes(M/E_{n-i+1})
  \xhookrightarrow{\ \tau_i\ }
  (M/E_{i-1})\otimes(M/E_{n-i+1})
\]
has kernel
\[
  (E_i/E_{i-1})\otimes E_{n-i+1}.
\]
The image of \(K_i\) in
\((E_i\otimes M)/(E_{i-1}\otimes M)\) is both contained in and
contains this kernel. Hence
\[
  K_i=K_{i-1}+E_i\otimes E_{n-i+1},
\]
which proves the lemma.

To conclude this paragraph, return to an arbitrary pseudocompact ring
\(k\).\footnote{\label{sga3:VIIB:nde:66}The editors added what
follows; the original treated only the case in which \(k\) is
Artinian.} Let
\((\mathbf H,\Delta,\varepsilon)\) be a flat
\(\mathbf O_k\)-coalgebra, put
\(\mathbf H^+=\ker(\varepsilon)\), and let
\[
  A=\Gamma^*(\mathbf H)
\]
be the dual profinite \(k\)-algebra. Put
\(X=\operatorname{Spf}(A)\), so that
\(\mathbf H=\mathbf H(X)\) \textup{(cf.\ \SGARefLink{sga3:VIIB:1.3.5}{1.3.5})}. Suppose that a
morphism of \(\mathbf O_k\)-coalgebras
\[
  \phi:\mathbf O_k\longrightarrow\mathbf H
\]
is given. It defines a continuous homomorphism of \(k\)-algebras
\(A\to k\), and hence a section \(\sigma\) of the structure morphism
\(X\to\operatorname{Spf}(k)\).

For every object \(B\) of \(\operatorname{Alf}/k\), put
\[
  \mathbf H_0(B)=\phi(B)=B\phi_B,
  \qquad
  \phi_B=\phi(1_B)\in\mathbf H(B).
\]
For \(n\geq1\), define \(\mathbf O_k\)-submodules
\(\mathbf H_n\) of \(\mathbf H\) by
\[
  \mathbf H_n(B)
  =
  \left\{
    u\in\mathbf H(B)\ \middle|\
    \Delta(u)-u\otimes\phi_B
    \in
    \mathbf H_{n-1}(B)\otimes\mathbf H^+(B)
  \right\}.
\]
This gives a filtration
\[
  \mathbf H_0\subset\mathbf H_1\subset\cdots
\]
of \(\mathbf H(X)\). The preceding result gives:

\medskip
\noindent\textbf{Proposition.}
\phantomsection
\label{sga3:VIIB:1.3.6.proposition}
{\itshape The section \(\sigma\) induces an isomorphism of the
underlying topological spaces if and only if
\[
  \mathbf H(X)=\bigcup_n\mathbf H_n.
\]}

